\bookchapter{Κεφάλαιο 64: Αποθήκευση δεδο$\mu$ένων και κλήσεις δικτύου}

\chapterintro{Μια εφαρ$\mu$ογή που $\mu$όνο «κοιτάζει» είναι νεκρή· εφαρ$\mu$ογή που αποθηκεύει δεδο$\mu$ένα και συνδέεται στο internet είναι ζωντανή. Σε αυτό το κεφάλαιο γνωρίζεις το «τρίο» της στρώσης δεδο$\mu$ένων iOS: UserDefaults ($\mu$ικρές ρυθ$\mu$ίσεις), SwiftData (δο$\mu$η$\mu$ένα δεδο$\mu$ένα), URLSession (επικοινωνία $\mu$ε τον διακο$\mu$ιστή), $\mu$αζί $\mu$ε το Codable για ανάλυση JSON και το $\mu$οντέρνο async/await. Μετά το κεφάλαιο αυτό, θα $\mu$πορείς να γράφεις πραγ$\mu$ατικές εφαρ$\mu$ογές που θυ$\mu$ούνται τις συνήθειες του χρήστη και ανανεώνουν περιεχό$\mu$ενο σε πραγ$\mu$ατικό χρόνο.}

\booksection{Α. Πώς διαλέγεις λύση αποθήκευσης: δες πρώτα τον πίνακα}

\begin{center}
\begin{tabularx}{\linewidth}{|X|X|X|X|}
\hline
\textbf{Λύση} & \textbf{Τι αποθηκεύει καλύτερα} & \textbf{Χαρακτηριστικά} & \textbf{Πού χρησι$\mu$οποιείται στο βιβλίο} \tabularnewline
\hline
UserDefaults & Σκόρπιες $\mu$ικρο-ρυθ$\mu$ίσεις (διακόπτες, προτι$\mu$ήσεις) & Ζεύγη κλειδί-τι$\mu$ής, απλό και ευθύ & Οι ρυθ$\mu$ίσεις στο Κεφ. 65, οι αγαπη$\mu$ένες πόλεις στο Κεφ. 66 \tabularnewline
SwiftData & Δο$\mu$η$\mu$ένα δεδο$\mu$ένα (συναλλαγές, καταγραφές) & Επίση$\mu$ο wrapper βάσης δεδο$\mu$ένων· δήλωση $\mu$ε $\mu$ία γρα$\mu$$\mu$ή @Model & Οι καταγραφές συναλλαγών στο Κεφ. 65 \tabularnewline
URLSession & (δεν αποθηκεύει, τα παίρνεις online) Δεδο$\mu$ένα διακο$\mu$ιστή & Ενσω$\mu$ατω$\mu$ένο στο σύστη$\mu$α framework δικτύου, χωρίς εξαρτήσεις από τρίτους & Τα δεδο$\mu$ένα καιρού στο Κεφ. 66 \tabularnewline
\hline
\end{tabularx}
\end{center}

\begin{center}\small\color{gray}Ο ρόλος του τρίο των δεδο$\mu$ένων\end{center}

Κράτα ένα $\mu$νη$\mu$όνιο: «$\mu$ικρές ρυθ$\mu$ίσεις στο UserDefaults, $\mu$αζικά δεδο$\mu$ένα στο SwiftData, τα δεδο$\mu$ένα απ' έξω περνούν από το URLSession». Το Core Data των παλιών οδηγών εξακολουθεί να δουλεύει (το SwiftData είναι η $\mu$οντέρνα εκδοχή του), ενώ το χειρόγραφο SQLite σχεδόν δεν το αγγίζεις — $\mu$άθε απευθείας το νέο, $\mu$ην πας παλιό δρό$\mu$ο.

\booksection{Β. UserDefaults: αποθηκεύεις $\mu$ικρές ρυθ$\mu$ίσεις}

\begin{lstlisting}[language=Swift]
// Γράψι$\mu$ο: αποθήκευση τι$\mu$ής
UserDefaults.standard.set(true, forKey: "isDarkMode")
UserDefaults.standard.set("Πεκίνο ", forKey: "favoriteCity")

// Διάβασ$\mu$α: ανάκτηση (χωρίς τι$\mu$ή, επιστρέφεται η προεπιλογή)
let isDark = UserDefaults.standard.bool(forKey: "isDarkMode")
let city = UserDefaults.standard.string(forKey: "favoriteCity") ?? "Πεκίνο"

// Διαγραφή: αφαίρεση κλειδιού
UserDefaults.standard.removeObject(forKey: "isDarkMode")
\end{lstlisting}

Πρόσεξε ότι ταιριάζει $\mu$όνο για σκόρπιες $\mu$ικρο-ρυθ$\mu$ίσεις — για δεκάδες ή εκατοντάδες εγγραφές άσε το στο SwiftData, και $\mu$ην χρησι$\mu$οποιείς το UserDefaults σαν βάση δεδο$\mu$ένων (καθυστερεί την εκκίνηση και δεν ευννοεί αναζητήσεις).

\booksection{Γ. SwiftData: δώσε στα δεδο$\mu$ένα «σπίτι αληθινό»}

Το SwiftData είναι η επίση$\mu$η λύση βάσης δεδο$\mu$ένων από το iOS 17, και έχει απλοποιήσει στο έπακρο τη «δη$\mu$ιουργία πινάκων και τις προσθήκες/διαγραφές/αλλαγές/αναζητήσεις» — $\mu$ε ένα annotation @Model, η κλάση δεδο$\mu$ένων γίνεται πίνακας βάσης:

\begin{lstlisting}[language=Swift]
import SwiftData

// 1. @Model: η κλάση γίνεται αυτό$\mu$ατα "πίνακας"
@Model
class Record {
    var amount: Double
    var category: String
    var note: String
    var date: Date

    init(amount: Double, category: String, note: String, date: Date) {
        self.amount = amount
        self.category = category
        self.note = note
        self.date = date
    }
}

// 2. Σύνδεση container στην είσοδο της εφαρ$\mu$ογής
@main
struct MyApp: App {
    var body: some Scene {
        WindowGroup {
            ContentView()
        }
        .modelContainer(for: Record.self)   // η βάση είναι έτοι$\mu$η
    }
}

// 3. Στο view ανάγνωση $\mu$ε @Query, εγγραφή $\mu$ε @Environment
struct RecordListView: View {
    @Query(sort: \Record.date, order: .reverse) var records: [Record]
    @Environment(\.modelContext) private var context

    var body: some View {
        List(records) { r in
            Text("\(r.category) \(r.amount) γιουάν")
        }
    }

    func add(amount: Double, category: String) {
        let r = Record(amount: amount, category: category, note: "", date: Date())
        context.insert(r)     // αποθήκευση
        try? context.save()   // εγγραφή στον δίσκο
    }
}
\end{lstlisting}

Αυτό το template «@Model + modelContainer + @Query» $\mu$πορείς να το $\mu$εταφέρεις αυτούσιο σε όποιο project θέλεις — απλά άλλαξε τα ονό$\mu$ατα των πεδίων. Μόλις αλλάξουν τα δεδο$\mu$ένα, το @Query ανανεώνει αυτό$\mu$ατα τη διεπαφή — $\mu$αγικό.

\booksection{Δ. Codable: «$\mu$εταφράζεις» το JSON}

Τα δεδο$\mu$ένα του διακο$\mu$ιστή είναι ως επί το πλείστον κεί$\mu$ενο JSON. Το protocol Codable της Swift κάνει τη $\mu$ετατροπή «JSON struct» να γίνεται αυτό$\mu$ατα — αν τα ονό$\mu$ατα των πεδίων ταιριάζουν, η $\mu$αγεία δουλεύει:

\begin{lstlisting}[language=Swift]
// JSON: {"name":"Πεκίνο","main":{"temp":26.5},"weather":[{"description":"Αίθριο"}]}
struct WeatherResponse: Codable {
    let name: String
    let main: Main
    let weather: [Weather]
}
struct Main: Codable { let temp: Double }
struct Weather: Codable { let description: String }

// Αλφαριθ$\mu$ητικό JSON  struct
let decoder = JSONDecoder()
let data = jsonString.data(using: .utf8)!
let result = try decoder.decode(WeatherResponse.self, from: data)
print(result.weather[0].description)   // "Αίθριο"
\end{lstlisting}

Σιδερένιος κανόνας: πρώτα κοιτά το πραγ$\mu$ατικό JSON που επιστρέφεται, $\mu$ετά όρισε τα πεδία. Πολλά bugs «αποτυχίας ανάλυσης» οφείλονται σε λάθος όνο$\mu$α πεδίου ή λάθος τύπο (ο διακο$\mu$ιστής στέλνει «26.5» σαν συ$\mu$βολοσειρά, εσύ το δήλωσες Double) — τύπωσε το αρχικό JSON και σύγκρινε, και θα το εντοπίσεις σε δέκα δευτερόλεπτα.

\booksection{Ε. URLSession + async/await: $\mu$ιλάς $\mu$ε τον διακο$\mu$ιστή}

Στο iOS η σύνδεση στο internet γίνεται $\mu$ε το ενσω$\mu$ατω$\mu$ένο URLSession, σε συνδυασ$\mu$ό $\mu$ε το async/await (ασύγχρονες συναρτήσεις) της Swift — το «περι$\mu$ένεις το internet» είναι τόσο φυσικό όσο το «περι$\mu$ένεις το courier»:

\begin{lstlisting}[language=Swift]
struct WeatherService {
    // Ασύγχρονη συνάρτηση: async ση$\mu$αίνει "πρέπει να περι$\mu$ένω λίγο", throws ση$\mu$αίνει "$\mu$πορεί να αποτύχει"
    func fetchWeather(city: String) async throws -> WeatherResponse {
        let url = URL(string:
            "https://api.example.com/weather?city=\(city)")!
        let (data, _) = try await URLSession.shared.data(from: url)
        return try JSONDecoder().decode(WeatherResponse.self, from: data)
    }
}

// Ση$\mu$είο κλήσης ($\mu$έσα σε .task):
.task {
    do {
        let weather = try await service.fetchWeather(city: "Πεκίνο")
        // Ενη$\mu$έρωση οθόνης……
    } catch {
        print("Σφάλ$\mu$α: \(error)")
    }
}
\end{lstlisting}

\textbf{Υπενθύ$\mu$ιση αποφυγής παγίδων: δύο σιδερένιοι κανόνες δικτύου}\par

(1) Οι κλήσεις δικτύου γίνονται πάντα $\mu$ε async/await, ποτέ δεν $\mu$πλοκάρεις το κύριο νή$\mu$α — η κλειδω$\mu$ένη διεπαφή είναι ε$\mu$πειρία επιπέδου crash. Η δο$\mu$η$\mu$ένη ταυτόχρονη εκτέλεση της Swift το κάνει φυσικό, αλλά $\mu$ην ξεχνάς το try await και το do/catch. (2) Το ATS απαγορεύει από προεπιλογή το HTTP σε καθαρό κεί$\mu$ενο: το iOS θέλει HTTPS. Για δοκι$\mu$ή τοπικού http:// service χρειάζεται προσωρινή εξαίρεση στο Info.plist (NSAppTransportSecurity), αλλά το τελικό προϊόν πρέπει να περνάει εξ ολοκλήρου από HTTPS — και η έγκριση το προσέχει.

\begin{tcbpractice}{Ας το δοκι$\mu$άσου$\mu$ε: φτιάξε $\mu$ια $\mu$ικρή εφαρ$\mu$ογή «φράση της η$\mu$έρας»}
\booksection{ΣΤ. Τα σφάλ$\mu$ατα δικτύου θέλουν «βαθ$\mu$ίδες»: $\mu$ην πετάς όλα τα λάθη στον χρήστη}

Χωρίς internet, λήξη χρόνου, 404, 500 διακο$\mu$ιστή... σε διαφορετικό λάθος, ο χρήστης πρέπει να δει διαφορετικό $\mu$ήνυ$\mu$α:

\begin{lstlisting}[language=Swift]
func friendlyMessage(_ error: Error) -> String {
    if let urlError = error as? URLError {
        switch urlError.code {
        case .notConnectedToInternet: return "Έπεσε η σύνδεση, να τη ελέγξεις;"
        case .timedOut:               return "Το αίτη$\mu$α έληξε, δοκι$\mu$ή αργότερα;"
        case .cannotFindHost:         return "Ο server δεν βρέθηκε"
        default:                      return "Πρόβλη$\mu$α δικτύου: \(urlError.code)"
        }
    }
    if let d = error as? DecodingError {
        return "Λάθος $\mu$ορφή δεδο$\mu$ένων (ίσως άλλαξε το API)"
    }
    return "Σφάλ$\mu$α: \(error.localizedDescription)"
}
\end{lstlisting}

Δείξε στον χρήστη «ανθρώπινη γλώσσα», και στο log άσε το «πρωτότυπο κεί$\mu$ενο» (print το πλήρες error). Μην πετάς στον χρήστη το αγγλικό $\mu$υστήριο του error.localizedDescription — η αίσθηση του προϊόντος κρύβεται σε αυτές τις λεπτο$\mu$έρειες.

Βρες ένα δωρεάν API φράσεων (π.χ. hitokoto ή παρό$\mu$οιο), και $\mu$ε URLSession + async/await φέρνε τη «φράση της η$\mu$έρας» και ε$\mu$φάνισέ τη στο κέντρο της οθόνης, προσθέτοντας ένα κου$\mu$πί «Άλλαξε φράση». Κατόπιν αποθήκευσε τις αγαπη$\mu$ένες φράσεις σε «Αγαπη$\mu$ένα» $\mu$ε SwiftData. Αυτή η άσκηση χρησι$\mu$οποιεί όλα τα UserDefaults/SwiftData/Codable/URLSession — αφού την τελειώσεις, τα Κεφάλαια 7 και 8 θα σου φανούν παιχνίδι.

\end{tcbpractice}

\begin{tcbsummary}{Σύνοψη κεφαλαίου}
\begin{itemize}
  \item Μνη$\mu$όνιο: $\mu$ικρές ρυθ$\mu$ίσεις στο UserDefaults, $\mu$αζικά δεδο$\mu$ένα στο SwiftData, τα εξωτερικά δεδο$\mu$ένα περνούν από το URLSession
  \item Τα τρία «χτυπή$\mu$ατα» του SwiftData: @Model δηλώνει τον πίνακα, modelContainer κρε$\mu$ά το container, @Query διαβάζει + context.insert γράφει
  \item Το Codable «$\mu$εταφράζει» αυτό$\mu$ατα το JSON: αν τα πεδία ταιριάζουν, η $\mu$αγεία δουλεύει
  \item URLSession + async/await: το internet σαν να περι$\mu$ένεις courier· το do/catch διαχειρίζεται τις αποτυχίες
  \item Δύο σιδερένιοι κανόνες: $\mu$ην $\mu$πλοκάρεις το κύριο νή$\mu$α· το τελικό προϊόν πρέπει να είναι HTTPS (ATS)
\end{itemize}

\end{tcbsummary}
